\section{OBJECTIVE AND SOURCE REVIEW}
\textbf{Erd\H{o}s \#644; independent attempt, 2026-09-23.}
Determine whether $f(k,7)=(3/4+o(1))k$, and whether $f(k,r)/k$ has a limit for every fixed $r\ge3$. Here $f(k,r)$ is the maximal transversal number for $k$-uniform families whose every subfamily of at most $r$ edges has a two-point transversal, using the supplied nonvacuous formulation. I read the full \texttt{644.tex}, including its complete-hypergraph lower bound and imported upper bound. I test how lower constructions transfer between uniformities and whether the usual vertex duplication improves the ratio.
\section{WORK: PADDING AND DUPLICATION HAVE DIFFERENT EFFECTS}
Given such a family $\mathcal H$, attach to each edge $E$ a distinct private point $z_E$ and replace it by $E\cup\{z_E\}$. This gives a $(k+1)$-uniform family $\mathcal H'$ with the same transversal number. An old transversal still works, giving $\tau(\mathcal H')\le\tau(\mathcal H)$. Conversely replace each private point of any transversal of $\mathcal H'$ by an arbitrary old point of its one associated edge. The resulting set has no greater size and meets every old edge. Hence equality holds. The local two-point property is preserved because each old local transversal still meets the enlarged edges. Thus $f(k+1,7)\ge f(k,7)$ in this formulation.
In contrast, replace every vertex $v$ by $t$ clones and each edge by the union of all clones of its vertices. The new rank is $tk$, but its transversal number remains $\tau(\mathcal H)$: choosing one clone per old transversal works, and projecting any new transversal to old vertices works in reverse. The local property again persists. Therefore this common duplication operation divides, rather than preserves, the ratio $\tau/k$ by $t$.
\section{ADVERSARIAL VERIFICATION AND REMAINING GAP}
Private points must be different for different edges; one shared new point would make the transversal number one. Duplicate selections may collapse on projection, which only helps the inequalities. Neither operation multiplies transversal number with rank, so neither promotes one isolated high-ratio example into an asymptotically equal ratio at all ranks. The source's known sandwich can be retained as external background, but no narrower interval or exact limiting constant is obtained. The at-most-seven formulation avoids vacuity for very small families.
Both transformations work verbatim for every fixed $r$, but monotonicity of $f(k,r)$ alone does not prove convergence of its ratio to $k$. No such limit is obtained here.
\section{SOURCES CHECKED}
Complete local source \texttt{gpt\_pro\_5.4/644.tex}; official indexed statement \url{https://www.erdosproblems.com/forum/thread/644}. Both transformation lemmas are proved directly, with no claim of novelty or use of the source's unproved external extremal bounds.
\section{FINAL}
\textbf{UNRESOLVED.} Rank padding is verified, and a tempting ratio-preserving duplication argument is ruled out.
\section{COMPLETION ESTIMATE}
Conservative subjective progress toward the exact extremal constant.
\textbf{COMPLETION: 8\%}
